\documentclass[a4paper]{exam}
\usepackage{amsmath,amssymb,titlesec}
\usepackage[a4paper,margin=22mm]{geometry}

\titleformat{\section}{\normalfont\Large\bfseries}{}{0pt}{}
\titleformat{\subsection}{\normalfont\large\bfseries}{}{0pt}{}
\titlespacing*{\section}{1em}{1.6\baselineskip}{0.6\baselineskip}
\titlespacing*{\subsection}{2em}{1.2\baselineskip}{0.4\baselineskip}

\setlength{\parindent}{0in}
\setlength{\headheight}{18pt}

\printanswers
\title{Week 8}
\author{Radhesh Goel}
\date{\today}

\begin{document}
\maketitle

\section*{Question 1.}

Rearrange the velocity equation into a root-finding problem and verify
that \(50\ \mathrm{kPa} \leq p \leq p_0\) contains a root.

\begin{solution}
  Subtracting the measured velocity gives
  \[
    \boxed{f(p)=\sqrt{2c_pT_0
      \left[1-\left(\frac{p}{p_0}\right)^{(\gamma-1)/\gamma}\right]}-700=0}
  \]
  With all pressures in pascals,
  \[
    f(p)=\sqrt{1809000\left[1-
      \left(\frac{p}{350000}\right)^{2/7}\right]}-700
  \]
  This function is continuous on \([50000,350000]\), since the square-root
  argument is non-negative throughout the interval. At its endpoints,
  \[
    \begin{aligned}
      f(50000)&\approx 178.359052\ \mathrm{m\,s^{-1}}>0,\\
      f(350000)&=-700\ \mathrm{m\,s^{-1}}<0.
    \end{aligned}
  \]
  Therefore, the Intermediate Value Theorem guarantees at least one root.
  The velocity decreases strictly as pressure increases, so this root is
  unique.
\end{solution}
\end{document}
